\section{Decoding Stochasticity versus Determinism}
\label{sec:sampling}
% Auto-generated by scripts/generate_tables.py -- do not edit.
\begin{table}[t]
\centering
\caption{Decoding stochasticity versus determinism. Sampled columns use temperature 0.7, top-$p$ 0.95 with three different seeds; the range across seeds is in percentage points. Greedy-seed rows change only the RNG seed under temperature 0: the committed runs are item-for-item identical to the control (0 discordant pairs out of 800), so greedy seed variation contributes exactly zero measured variance in this harness.}
\label{tab:sampling}
\small
\begin{tabular}{lcccccc}
\toprule
Model & Greedy (T=0) & T=0.7, s0 & T=0.7, s1 & T=0.7, s2 & Range (pp) & Greedy-seed no-op \\
\midrule
Qwen2.5-3B & 64.4 & 64.2 & 65.0 & 63.0 & 2.00 & yes (0/800) \\
SmolLM2-1.7B & 39.8 & 36.2 & 37.6 & 39.0 & 2.75 & yes (0/800) \\
\bottomrule
\end{tabular}
\end{table}


Two very different phenomena hide behind the word ``seed.'' We measure both
in the same harness (\cref{tab:sampling}, \cref{fig:sampling}). For the
consolidated view of all configurations including these results, see
\cref{tab:master}.

\paragraph{Greedy seed changes are exact no-ops.}
Under temperature 0, changing the RNG seed (0$\to$1, 0$\to$2) reproduces the
control run item-for-item in all six greedy-seed contrasts (0 discordant
pairs out of 800 in each). In this harness, greedy decoding is deterministic
across seeds, so greedy seed variation contributes exactly zero measured
variance. Evaluations that report ``different-seed greedy runs'' as
independent replicates should not expect them to vary under deterministic
decoding; stochasticity must be introduced through sampling.

\paragraph{Temperature-0.7 sampling produces churn with small drift.}
At temperature 0.7, top-$p$ 0.95, the three seed replicates score:
Qwen2.5-3B 64.25\%, 65.00\%, 63.00\% (range 2.00\,\pp); SmolLM2-1.7B 36.25\%,
37.63\%, 39.00\% (range 2.75\,\pp). The seed-to-seed spread is thus
comparable to the quantization and backend effects, and small compared to
prompt effects. The per-item picture is very different: 9.9--24.1\% of
items change correctness state between any sampled run and the greedy
control (the Qwen range is 9.9--11.1\%; the SmolLM2 range is 21.0--24.1\%). Sampling noise is therefore far from a uniform small perturbation
at the item level; it is a high-churn, low-drift process whose net effect
partially cancels. A single sampled run can easily land $\pm$2\,\pp{} away
from another seed's score purely by chance; on a leaderboard where the
top-two gap is of this size (here 2.62\,pp{} with overlapping intervals),
one sampling seed is not an adequate summary
of a model.

\begin{figure}[t]
\centering
\includegraphics[width=0.66\linewidth]{sampling_variance.pdf}
\caption{Greedy decoding is seed-invariant (three identical points per
model; 0 discordant items in 800), while temperature-0.7 sampling (three
seeds, Wilson 95\% intervals) produces seed-to-seed ranges of 2.0--2.8\,\pp{}
on top of a systematic drift relative to greedy.}
\label{fig:sampling}
\end{figure}

\section{Ranking Stability}
\label{sec:ranking}
% Generated by scripts/audited_analysis.py
\begin{table}[tbp]
\centering\small
\caption{Ranks ordered Smol/Qwen/Phi, Kendall tau-b against control, and pairwise reversals. Ties would receive shared ranks; none occur in these columns.}\label{tab:ranking}
\begin{tabular}{lrrr}\toprule
Configuration & Ranks & $\tau_b$ & Reversals \\
\midrule
control & 3/2/1 & 1.000 & 0 \\
concise & 3/2/1 & 1.000 & 0 \\
5shot & 3/1/2 & 0.333 & 1 \\
vllm & 3/2/1 & 1.000 & 0 \\
seed 1 & 3/2/1 & 1.000 & 0 \\
seed 2 & 3/2/1 & 1.000 & 0 \\
int8 & 3/2/1 & 1.000 & 0 \\
int4 & 3/2/1 & 1.000 & 0 \\
\bottomrule\end{tabular}
\end{table}

% Generated by scripts/audited_analysis.py
\begin{table}[tbp]
\centering\small
\caption{Qwen minus Phi: paired model comparison. Holm correction covers the separate family of shared-configuration comparisons. Positive values favor Qwen. Conditional, post-hoc analysis; no claim of equivalence from non-significance.}\label{tab:h2h}
\begin{tabular}{lrrr}\toprule
Configuration & Difference (pp) & 95\% interval & Holm $p$ \\
\midrule
control & -2.625 & [-5.62, 0.38] & 0.47 \\
concise & -7.625 & [-10.62, -4.62] & $<0.001$ \\
5shot & +12.250 & [9.00, 15.50] & $<0.001$ \\
vllm & -0.375 & [-3.38, 2.50] & 0.87 \\
seed 1 & -2.625 & [-5.62, 0.38] & 0.47 \\
seed 2 & -2.625 & [-5.62, 0.38] & 0.47 \\
int8 & -2.500 & [-5.50, 0.62] & 0.47 \\
int4 & -4.250 & [-7.38, -1.25] & 0.0559 \\
\bottomrule\end{tabular}
\end{table}


\subsection{Scores move; orderings mostly hold}
\cref{tab:ranking} arranges the eight configurations measured for all three
fully-evaluated models and computes, per configuration, the competition
ranks and Kendall's $\taub$ against the control-config ranking. Under every
measured factor except the five-shot prompt, the ranking is identical to
the control ($\taub = 1.0$, zero reversals): Phi-3.5-mini first, Qwen2.5-3B
second, SmolLM2-1.7B third. The five-shot prompt produces the single
measured ranking change: $\taub = 0.33$ with one pairwise reversal, as
Qwen2.5-3B (68.63\%) overtakes Phi-3.5-mini (56.37\%).

\subsection{The reversal is interval-supported, and the control ordering is a tie}
\cref{tab:h2h} examines the Phi/Qwen head-to-head. Under control (and under
seeds, backend, and quantization) Phi-3.5-mini leads Qwen2.5-3B by
$+0.38$ to $+4.25$\pp, but in all of those configurations the two Wilson
intervals overlap: the control ordering is a point-estimate ordering on top
of a statistical tie. Under the five-shot prompt the difference is
$-12.26$\pp{} \emph{with disjoint intervals}---the only configuration in
which either model statistically separates the other. We therefore state
the finding conservatively: the measured configuration changes preserve a
non-separated Phi$\geq$Qwen ordering, while one legitimate prompt change
produces a statistically separated ordering in the opposite direction.

Pairwise win rates over the eight shared configurations quantify the same
asymmetry (\cref{fig:winprob}): Phi-3.5-mini beats Qwen2.5-3B in 87.5\% of
configurations (7 of 8), Qwen2.5-3B and Phi-3.5-mini each beat SmolLM2-1.7B
in 100\%. Mean competition ranks across the eight configurations are
1.13 (Phi), 1.88 (Qwen), 3.00 (Smol); P(first place) is 87.5\% / 12.5\% /
0\%. The seeded bootstrap over configurations (2{,}000 replicates, resampling
the eight configurations with replacement) yields mean $\taub=0.84$
against the mean-accuracy reference ranking and a 24.0\% probability that a
resampled matrix produces at least one pairwise reversal relative to that
reference---i.e., the top-two order is stable in the point-estimate
sequence but only weakly separated statistically, consistent with the
interval analysis.

\subsection{Score instability vs.\ ranking instability}
These are distinct quantities and both matter. The concise prompt moves
scores enormously ($-8.1$ to $-16.5$\pp) yet changes no ordering
($\taub=1.0$): every model degrades, and the relative ordering survives.
The five-shot prompt moves scores less dramatically for two of three models
and yet changes the \emph{decision} a leaderboard user would make, because
the two strongest models swap places. A small score change is consequential
when it crosses a ranking boundary between nearly-equal models; a large
score change can be inconsequential for selection when it is uniform. This
is the central practical lesson of the study for this model pair and task
suite.

\begin{figure}[t]
\centering
\includegraphics[width=0.5\linewidth]{rank_heatmap.pdf}
\caption{Model rank (1 = best) per evaluation configuration. The outlined
five-shot column is the only ranking change: Qwen2.5-3B and Phi-3.5-mini
swap first and second place ($\taub=0.33$ vs.\ the control ranking, one
pairwise reversal); SmolLM2-1.7B is third throughout.}
\label{fig:rankheat}
\end{figure}

\begin{figure}[t]
\centering
\includegraphics[width=0.42\linewidth]{win_probability.pdf}
\caption{Pairwise win rates over the eight shared configurations:
fraction of configurations in which the row model's point estimate exceeds
the column model's. Phi-3.5-mini leads Qwen2.5-3B in 7 of 8 configurations;
both lead SmolLM2-1.7B in all 8.}
\label{fig:winprob}
\end{figure}